\documentclass[../../../include/open-logic-section]{subfiles}

\begin{document}

\olfileid{sth}{cardinals}{classing}
\olsection{સાન્ત, \usetoken{S}{enumerable}, \usetoken{S}{nonenumerable}}

ગણાંકોનો પરિચય થયા પછી તેમના જુદા પ્રકારો વિશે થોડું
વિચારવું યોગ્ય છે; ખાસ કરીને સાન્ત, !!{enumerable}
અને !!{nonenumerable} ગણાંકો વિશે.

આપણા પહેલા બે પરિણામોથી સાન્ત ગણાંકો બરાબર સાન્ત
ક્રમવાચક સંખ્યાઓ જ હશે. તેમને આપણે અગાઉ
\olref[z][infinity-again]{defnomega}માં આપણી
\emph{પ્રાકૃતિક સંખ્યાઓ} તરીકે વ્યાખ્યાયિત કરી હતી:

\begin{prop}\ollabel{finitecardisoequal}
ધારો કે $n, m \in \omega$. ત્યારે $n = m$ જો અને માત્ર જો
$\cardeq{n}{m}$.
\end{prop}

\begin{proof}
\emph{ડાબેથી જમણે} સ્પષ્ટ છે. \emph{જમણેથી ડાબે}
માટે ધારો કે $\cardeq{n}{m}$ છતાં $n \neq m$.
ત્રિવિભાજનથી કાં તો $n \in m$ અથવા $m \in n$;
સામાન્યતા ગુમાવ્યા વિના $n \in m$ ધારો.
ત્યારે $n \subsetneq m$ અને
!!a{bijection} $f \colon m \to n$ છે, તેથી $m$
ડેડેકિન્ડ-અનંત છે. આ
\olref[z][infinity-again]{naturalnumbersarentinfinite}
સાથે વિરોધાભાસ છે.
\end{proof}

\begin{cor}\ollabel{naturalsarecardinals}
જો $n \in \omega$, તો $n$ ગણાંક છે.
\end{cor}

\begin{proof}
તરત જ મળે છે.
\end{proof}
\noindent
આમાંથી એ પણ મળે છે કે ગણાંકને ``સાન્ત'' અથવા
``અનંત'' કહેવાના ઘણા યોગ્ય લાગે એવા ખ્યાલો એકસરખા છે:
\begin{thm}\ollabel{generalinfinitycharacter}કોઈ પણ ગણ $A$ માટે નીચેની શરતો સમકક્ષ છે:
\begin{enumerate}
	\item\ollabel{card:notinomega} $\card{A} \notin \omega$, એટલે કે
	$\card{A}$ પ્રાકૃતિક સંખ્યા નથી;
	\sourcecorrection{OLMD-135}{મૂળમાં અહીં ``$A$ પ્રાકૃતિક સંખ્યા નથી'' લખાયું છે, પરંતુ
	$A$ મનસ્વી ગણ છે અને સાન્ત હોવા છતાં પ્રાકૃતિક સંખ્યા ન હોઈ શકે.
	શરત $\card{A} \notin \omega$ને સમકક્ષ અર્થમાં
	``$\card{A}$ પ્રાકૃતિક સંખ્યા નથી'' વાંચવું જોઈએ.}
	\item\ollabel{card:omegaplus} $\omega \leq \card{A}$;
	\item\ollabel{card:infinite} $A$ ડેડેકિન્ડ-અનંત છે.
\end{enumerate}
\end{thm}

\begin{proof}
\olref[ord-arithmetic][using-addition]{ordinfinitycharacter},
\olref[cardinals][cardsasords]{lem:CardinalsBehaveRight},
અને \olref{naturalsarecardinals}થી મળે છે.
\end{proof}

આથી \olref[sfr][][]{part}માં થોડા અનૌપચારિક રીતે
વાપરેલા કેટલાક ખ્યાલોની નીચેની \emph{વ્યાખ્યા} આપી શકીએ:

\begin{defn}\ollabel{defnfinite}
$A$ને \emph{સાન્ત} કહીએ જો અને માત્ર જો
$\card{A}$ પ્રાકૃતિક સંખ્યા હોય, એટલે કે
$\card{A} \in \omega$. અન્યથા $A$ને
\emph{અનંત} કહીએ.
\end{defn}
\noindent
પરંતુ નોંધો કે આ વ્યાખ્યા $\ZFC$ની પૃષ્ઠભૂમિમાં
આપેલી છે. છેવટે દરેક ગણનો ગણાંક હોય તેની ખાતરી
માટે આપણને સુક્રમિતતા જરૂરી હતી. ખરેખર સુક્રમિતતા
વિના એવો ગણ હોઈ શકે જે ન તો સાન્ત હોય, ન તો
ડેડેકિન્ડ-અનંત. આ પ્રકારના પ્રશ્ન પર
\olref[choice][]{chap}માં ફરી આવીશું. હમણાં
સુક્રમિતતા વાપરવાનું ચાલુ રાખીએ.

હવે સાન્ત ગણાંકોથી અનંત ગણાંકો તરફ વળીએ.
પહેલી બે સાદી બાબતો આ છે:

\begin{cor}\ollabel{omegaisacardinal}
$\omega$ લઘુતમ અનંત ગણાંક છે.
\end{cor}

\begin{proof}
$\omega$ ગણાંક છે, કારણ કે $\omega$ ડેડેકિન્ડ-અનંત
છે અને કોઈ $n \in \omega$ માટે
$\cardeq{\omega}{n}$ હોય, તો $n$ પણ ડેડેકિન્ડ-અનંત
થાય; આ \olref[z][infinity-again]{naturalnumbersarentinfinite}
સાથે વિરોધાભાસ છે. હવે વ્યાખ્યાથી $\omega$
લઘુતમ અનંત ગણાંક છે.
\end{proof}

\begin{cor}
દરેક અનંત ગણાંક સીમા ક્રમવાચક સંખ્યા છે.
\end{cor}

\begin{proof}
ધારો કે $\alpha$ અનંત અનુગામી ક્રમવાચક સંખ્યા છે,
એટલે કોઈ $\beta$ માટે $\alpha = \beta
\ordplus 1$. $\beta$ પણ અનંત છે: મૂળમાં અહીં
\olref{finitecardisoequal}નો ઉલ્લેખ છે, પરંતુ
પૂર્વવર્તી ક્રમવાચક સંખ્યા સાન્ત હોય, તો તેનો
અનુગામી પણ સાન્ત થાય. તેથી
\olref[ord-arithmetic][using-addition]{ordinfinitycharacter}
મુજબ $\cardeq{\beta}{\beta \ordplus 1}$.
હવે \olref[cardinals][cardsasords]{lem:CardinalsBehaveRight}
મુજબ $\card{\beta} = \card{\beta\ordplus 1} =
\card{\alpha}$, તેથી $\alpha \neq \card{\alpha}$.
\sourcecorrection{OLMD-136}{મૂળે $\beta$ અનંત છે તે માટે
\olref{finitecardisoequal}નો સંદર્ભ આપ્યો છે; તે
પરિણામ ફક્ત બે સાન્ત ક્રમવાચક સંખ્યાઓની સમકદતાની
ચર્ચા કરે છે. અહીં જરૂરી કારણ એ છે કે સાન્ત
$\beta$નો અનુગામી પણ સાન્ત થાય, જે અનંત
$\alpha$ને વિરોધે છે.}
\end{proof}

અગાઉ \olref[sfr][siz][enm-alt]{defn:enumerable}થી જ
આપણે !!{enumerable} અને !!{nonenumerable} અનંત
ગણોમાં ભેદ પાડ્યો હતો. એ વ્યાખ્યાથી સ્વાભાવિક રીતે
આ મળે છે:

\begin{prop}
$A$ એ !!{enumerable} હોય જો અને માત્ર જો
$\card{A} \leq \omega$, અને $A$ એ
!!{nonenumerable} હોય જો અને માત્ર જો
$\omega < \card{A}$.
\end{prop}

\begin{proof}
ત્રિવિભાજનથી બંને દાવા સમકક્ષ છે; તેથી
$A$ એ !!{enumerable} હોય જો અને માત્ર જો
$\card{A} \leq \omega$ તે સાબિત કરવું પૂરતું છે.
\emph{જમણેથી ડાબે}: $\card{A} \leq \omega$ હોય, તો
\olref[cardinals][cardsasords]{lem:CardinalsBehaveRight}
અને \olref{omegaisacardinal} મુજબ
$\cardle{A}{\omega}$.
\emph{ડાબેથી જમણે}: ધારો કે $A$ એ
!!{enumerable} છે; તો
\olref[sfr][siz][enm-alt]{defn:enumerable} મુજબ
ત્રણ સંભવિત કિસ્સા છે:
\begin{enumerate}
	\item જો $A = \emptyset$, તો
	\olref{naturalsarecardinals} અને
	\olref[cardinals][cardsasords]{lem:CardinalsBehaveRight}
	મુજબ $\card{A} = 0 \in \omega$.
	\item જો $\cardeq{n}{A}$, તો
	\olref{naturalsarecardinals} અને
	\olref[cardinals][cardsasords]{lem:CardinalsBehaveRight}
	મુજબ $\card{A} = n \in \omega$.
	\item જો $\cardeq{\omega}{A}$, તો
	\olref{omegaisacardinal} મુજબ $\card{A} = \omega$.
\end{enumerate}
તેથી બધા કિસ્સામાં $\card{A} \leq \omega$.
\end{proof}
\noindent
ખરેખર $\omega$ને ખાસ સ્થાન છે. ગણનીય ક્રમવાચક
સંખ્યાઓ ઘણી હોવા છતાં:

\begin{cor}
$\omega$ એકમાત્ર !!{enumerable} અનંત ગણાંક છે.
\end{cor}

\begin{proof}
ધારો કે $\cardfont{a}$ એ !!a{enumerable} અનંત ગણાંક
છે. $\cardfont{a}$ અનંત હોવાથી
$\omega \leq \cardfont{a}$. $\cardfont{a}$ એ
!!a{enumerable} ગણાંક હોવાથી
$\cardfont{a} = \card{\cardfont{a}} \leq \omega$.
તેથી ત્રિવિભાજનથી $\cardfont{a} = \omega$.
\end{proof}

અલબત્ત, અનંત ઘણા ગણાંકો છે. તેથી પૂછી શકાય:
\emph{ગણાંકો કેટલા છે?} આગળનાં પરિણામો બતાવે છે
કે કદાચ આ પ્રશ્ન ફરી વિચારવો જોઈએ.

\begin{prop}\ollabel{unioncardinalscardinal}
જો $X$નો દરેક ઘટક ગણાંક હોય, તો $\bigcup X$
પણ ગણાંક છે.
\end{prop}

\begin{proof}
$\bigcup X$ ક્રમવાચક સંખ્યા છે તે તપાસવું સરળ છે.
ધારો કે $\alpha \in \bigcup X$ ક્રમવાચક સંખ્યા છે;
તો કોઈ ગણાંક $\cardfont{b}$ માટે
$\alpha \in \cardfont{b} \in X$. $\cardfont{b}$
ગણાંક હોવાથી $\cardless{\alpha}{\cardfont{b}}$.
$\cardfont{b} \subseteq \bigcup X$ હોવાથી
$\cardle{\cardfont{b}}{\bigcup X}$, તેથી
$\cardneq{\alpha}{\bigcup X}$. સામાન્યકરણથી
$\bigcup X$ ગણાંક છે.
\end{proof}

\begin{thm}\ollabel{lem:NoLargestCardinal}
સૌથી મોટો ગણાંક નથી.
\end{thm}

\begin{proof}
કોઈ પણ ગણાંક $\cardfont{a}$ માટે કૅન્ટરનું પ્રમેય
(\olref[sfr][siz][car]{thm:cantor}) અને
\olref[cardinals][cardsasords]{lem:CardinalsExist}
મુજબ $\cardfont{a} < \card{\Pow{\cardfont{a}}}$.
\end{proof}

\begin{thm}
બધા ગણાંકોનો ગણ અસ્તિત્વમાં નથી.
\end{thm}

\begin{proof}
વિરોધાભાસ માટે ધારો કે $C = \Setabs{\cardfont{a}}{\cardfont{a} \text{
is a cardinal}}$. હવે \olref{unioncardinalscardinal}
મુજબ $\bigcup C$ ગણાંક છે; તેથી
\olref{lem:NoLargestCardinal} મુજબ કોઈ ગણાંક
$\cardfont{b} > \bigcup C$ છે. વ્યાખ્યાથી
$\cardfont{b} \in C$, એટલે
$\cardfont{b} \subseteq \bigcup{C}$, તેથી
$\cardfont{b} \leq \bigcup C$; આ વિરોધાભાસ છે.
\end{proof}

આ પરિણામને રસેલના વિરોધાભાસ અને બુરાલી-ફોર્ટીના
વિરોધાભાસ બંને સાથે સરખાવો.

\end{document}
